% Propietario conservado: /Users/ruben/Documents/ChatGPT/jueces y controles/REVISION_ARGUMENTAL_APERTURAS_CIERRES_20260915/II/manuscript_es/sections/07b_geometria_elipse.tex
% Adición separada: no sustituye 07_elipse.tex.
% Procedencia: integral 2249, 52.4.1--2, 52.6.2--3, 53.28 y 90.4.6.
% Recuperación de resultados y explicitación geométrica; certificado local
% verificar_elipse_II.py. No usa un radio físico como entrada generadora.

\subsection{Semiejes, focos y radios interiores}
\label{ii:sec:ii-elipse-geometria}

La figura se obtiene de las salidas angulares y de acción de las secciones
precedentes. APP conserva las hojas y sus residuos y cocientes; TRIT fija
la orientación; TPK transporta frontera, acarreo y memoria hasta las
publicaciones del mismo estado enriquecido y de la estructura discreta
conjunta del continuo. Aquí se recibe el par \((A,C^*)\) y la sección
positiva \(\hbar\) después de esa construcción. La geometría explicita
qué información conservan sus lectores, sin usar una longitud observada
para seleccionar el estado fuente. La sección~\ref{ii:sec:ii-elipse-geometria} desarrolla esa realización a partir de los lectores ya generados, y la construcción angular se conserva en la sección~\ref{v:v:ant:sec:ii-angulos}.

Sea \(\xi=C^*/A\in(-1,1)\), con \(A\ne0\) y ambos ángulos en la
misma unidad. Escribimos \(\eta=\operatorname{artanh}\xi\) y
\[
 a_S=\sqrt{2\hbar}\,e^{\eta/2},\qquad
 b_S=\sqrt{2\hbar}\,e^{-\eta/2},\qquad
 \gamma(\theta)=(a_S\cos\theta,b_S\sin\theta).
\]
Los semiejes están etiquetados por sus coordenadas: si \(\eta<0\),
el eje vertical es el mayor. Ambos pertenecen a la recta de raíz de
acción de la carta equilibrada, no a la recta de longitudes espaciales.

\begin{proposition}[Geometría focal e interior radial]
\label{ii:prop:ii-elipse-focos}
Sean \(a=\max(a_S,b_S)\), \(b=\min(a_S,b_S)\). La semidistancia
focal y la excentricidad son
\begin{equation}
 f_S=\sqrt{a^2-b^2}
 =\sqrt{4\hbar\sinh|\eta|},\qquad
 e_f=\frac{f_S}{a}=\sqrt{1-e^{-2|\eta|}}.
 \label{ii:eq:ii-elipse-focal}
\end{equation}
Para \(\eta>0\), los focos son \((\pm f_S,0)\); para
\(\eta<0\), son \((0,\pm f_S)\). En \(\eta=0\) coinciden con
el centro. El radio de frontera en dirección polar \(\phi\) es
\begin{equation}
 \rho_{\partial}(\phi)
 =\left(\frac{\cos^2\phi}{a_S^2}
        +\frac{\sin^2\phi}{b_S^2}\right)^{-1/2}.
 \label{ii:eq:ii-elipse-radio-polar}
\end{equation}
El segmento interior en esa dirección consta de
\(r(\cos\phi,\sin\phi)\), con \(0\le r<\rho_{\partial}(\phi)\).
En particular, \(b\le\rho_{\partial}(\phi)\le a\).
\end{proposition}

\begin{proof}
Los cuadrados de los semiejes mayor y menor son
\(2\hbar e^{|\eta|}\) y \(2\hbar e^{-|\eta|}\). Su diferencia
da \eqref{ii:eq:ii-elipse-focal}. Para verificar el foco, supongamos
\(\eta>0\). En \(X=\gamma(\theta)\), las distancias a
\(F_+=(f_S,0)\) y \(F_-=(-f_S,0)\) satisfacen
\[
 |X-F_+|^2=(a_S-f_S\cos\theta)^2,\qquad
 |X-F_-|^2=(a_S+f_S\cos\theta)^2.
\]
Como \(a_S>f_S\), sus raíces positivas suman \(2a_S\).
Para \(\eta<0\), intercambiar las coordenadas da la misma prueba.
En el caso circular, \(f_S=0\). Finalmente, sustituir
\((Q,P)=r(\cos\phi,\sin\phi)\) en
\(Q^2/a_S^2+P^2/b_S^2=1\) da
\eqref{ii:eq:ii-elipse-radio-polar}. El coeficiente de \(r^2\) está
entre \(a^{-2}\) y \(b^{-2}\), lo que prueba las cotas y la
caracterización del interior.
\end{proof}

\subsection{Fase paramétrica, ángulo polar y área de acción}
\label{ii:sec:ii-elipse-fase-area}

\begin{proposition}[Ley areal y dos coordenadas angulares]
\label{ii:prop:ii-elipse-area}
La fase \(\theta\) de \(\gamma\) y su ángulo polar continuo
\(\phi\), levantados desde cero, satisfacen
\begin{equation}
 \phi(\theta)=\arg(a_S\cos\theta+i b_S\sin\theta),\qquad
 \frac{d\phi}{d\theta}
 =\frac{a_Sb_S}{a_S^2\cos^2\theta+b_S^2\sin^2\theta}>0.
 \label{ii:eq:ii-elipse-fases}
\end{equation}
En una carta donde ambas tangentes están definidas,
\(\tan\phi=(b_S/a_S)\tan\theta\). El área orientada barrida es
\begin{equation}
 \mathcal A(\theta)=\frac12\int_0^\theta(QP'-PQ')\,dt
 =\hbar\theta,\qquad
 \mathcal A(1)=\hbar,\qquad \mathcal A(2\pi)=h.
 \label{ii:eq:ii-elipse-area-fase}
\end{equation}
\end{proposition}

\begin{proof}
La derivada del argumento de un vector no nulo es
\((QP'-PQ')/(Q^2+P^2)\). Aquí el numerador es
\(a_Sb_S=2\hbar\), lo que prueba la derivada y la monotonía.
El cociente \(P/Q\) da la relación de tangentes; el levantamiento
continuo fija los cuadrantes y el número de vueltas. La integral areal
vale \(a_Sb_S\theta/2=\hbar\theta\), y
\(h=2\pi\hbar\) da la última identidad.
\end{proof}

La fase paramétrica no coincide en general con el ángulo polar. Un radián
de fase barre \(\hbar\), pero no se identifica por ello con un arco
circular ni con un sector de ángulo polar unitario. En coordenadas polares,
la misma ley se escribe
\(d\mathcal A=\rho_{\partial}(\phi)^2d\phi/2\).
Con la orientación de \(\gamma\) utilizada aquí,
\(\oint P\,dQ=-h\): el área positiva es \(h\), mientras la
integral de esa uno-forma conserva su signo. Esto concuerda con la
distinción entre los transportes simpléctico y antisimpléctico de la
sección anterior.

\begin{figure}[htbp]
 \centering
 
% BEGIN OWNER /Users/ruben/Documents/ChatGPT/jueces y controles/REVISION_ARGUMENTAL_APERTURAS_CIERRES_20260915/II/manuscript_es/figures/elipse_accion_II.tex
% Geometría calculada, no coordenadas focales elegidas a ojo.
% Caso algebraico de ilustración xi=1/3; no evaluación de las constantes HMT.
\begingroup
\pgfmathsetmacro{\elXi}{1/3}
\pgfmathsetmacro{\elEta}{0.5*ln((1+\elXi)/(1-\elXi))}
\pgfmathsetmacro{\elA}{exp(\elEta/2)}
\pgfmathsetmacro{\elB}{exp(-\elEta/2)}
\pgfmathsetmacro{\elF}{sqrt(\elA*\elA-\elB*\elB)}
\pgfmathsetmacro{\elTheta}{30}
\pgfmathsetmacro{\elX}{\elA*cos(\elTheta)}
\pgfmathsetmacro{\elY}{\elB*sin(\elTheta)}
\pgfmathsetmacro{\elPolar}{atan(\elY/\elX)}
\begin{tikzpicture}[x=2.18cm,y=2.18cm,font=\small,
 every node/.style={inner sep=2pt},>=Stealth]
 \begin{scope}
  \node[font=\bfseries] at (0,1.38) {Forma y orientación};
  \fill[accent!12] (0,0)--(\elA,0)
    plot[domain=0:30,samples=50] ({\elA*cos(\x)},{\elB*sin(\x)})--cycle;
  \draw[ink,thick] (0,0) ellipse[x radius=\elA,y radius=\elB];
  \draw[->,black!55] (-1.42,0)--(1.48,0) node[right]{\(u\)};
  \draw[->,black!55] (0,-1.05)--(0,1.08) node[above]{\(v\)};
  \draw[accent,thick,-{Stealth}] (\elA,0)
    plot[domain=0:30,samples=50] ({\elA*cos(\x)},{\elB*sin(\x)});
  \draw[accent,thick] (0,0)--(\elX,\elY);
  \node[below left,font=\scriptsize] at (0,0) {\(O\)};
  \draw[accent] (.35,0) arc[start angle=0,end angle=\elPolar,radius=.35];
  \node[accent] at (.47,.065) {\(\phi\)};
  \fill[accent] (\elX,\elY) circle[radius=1.4pt];
  \node[above right] at (\elX,\elY) {\(X(\theta)\)};
  \node[accent,align=center,font=\scriptsize] at (1.26,.18) {\(\theta=\pi/6\)};
  \draw[dashed,black!45] (-\elF,0)--(\elX,\elY);
  \draw[dashed,black!45] (\elF,0)--(\elX,\elY);
  \fill[ink] (-\elF,0) circle[radius=1.5pt] node[below]{\(F_-\)};
  \fill[ink] (\elF,0) circle[radius=1.5pt] node[below]{\(F_+\)};
  \draw[<->,black!60] (-\elA,-1.04)--(0,-1.04)
    node[midway,below,font=\scriptsize]{\(a_S/\sqrt{2\hbar}\)};
  \draw[<->,black!60] (-1.32,0)--(-1.32,\elB)
    node[midway,left,font=\scriptsize]{\(\frac{b_S}{\sqrt{2\hbar}}\)};
  \node[align=center,font=\scriptsize] at (0,-1.55)
    {\(\rho_{\partial}(\phi)/\sqrt{2\hbar}=|OX|\)\\
     \(f_S^2=a_S^2-b_S^2\), \quad área total \(h\)};
 \end{scope}
 \begin{scope}[shift={(3.65,0)}]
  \node[font=\bfseries] at (0,1.38) {Forma recíproca};
  \draw[densely dashed,black!25] (0,0) ellipse[x radius=\elA,y radius=\elB];
  \draw[ink,thick] (0,0) ellipse[x radius=\elB,y radius=\elA];
  \draw[->,black!55] (-1.28,0)--(1.38,0) node[right]{\(u\)};
  \draw[->,black!55] (0,-1.30)--(0,1.26) node[right]{\(v\)};
  \fill[ink] (0,\elF) circle[radius=1.5pt] node[right]{\(F'_+\)};
  \fill[ink] (0,-\elF) circle[radius=1.5pt] node[right]{\(F'_-\)};
  \node[align=center,font=\scriptsize] at (0,-1.55)
    {\(\eta\mapsto-\eta\), \quad \(a_S\leftrightarrow b_S\)\\
     \(R_{\rm el}\mapsto R_{\rm el}^{-1}\), \quad \(a_Sb_S=2\hbar\)};
 \end{scope}
 \draw[->,accent,thick] (1.60,.80)--(2.17,.80)
   node[midway,above,font=\scriptsize]{intercambio};
\end{tikzpicture}
\endgroup

% END OWNER

 \caption{Semiejes, focos y reciprocidad de la elipse de acción en las
 coordenadas normalizadas \(u=Q/\sqrt{2\hbar}\),
 \(v=P/\sqrt{2\hbar}\). El sector marcado tiene fase paramétrica
 \(\theta=\pi/6\), distinta de su ángulo polar \(\phi\); su área
 de acción es \(\hbar\pi/6\). La figura utiliza el caso algebraico
 \(\xi=1/3\) para hacer legibles las relaciones: no atribuye ese valor
 a la salida angular HMT. Los focos se calculan desde los semiejes.
 A la derecha, el intercambio de ejes conserva el área total \(h\)
 e invierte el radio modular adimensional.}
 \label{ii:fig:ii-elipse-calculada}
\end{figure}

Para el caso de control de la figura, \(\xi=1/3\), las coordenadas
normalizadas son exactamente
\[
 \eta=\tfrac12\log2,\quad
 \frac{a_S}{\sqrt{2\hbar}}=2^{1/4}=1.1892071150\ldots,\quad
 \frac{b_S}{\sqrt{2\hbar}}=
 \frac{f_S}{\sqrt{2\hbar}}=R_{\rm el}=2^{-1/4}=0.8408964153\ldots.
\]
Son números de regresión del dibujo, no valores atribuidos al generador
angular. El certificado local comprueba además la forma circular y las
orientaciones \(\xi=\pm1/3,\pm3/5\).

\subsection{Reciprocidad de forma y distancia interior}

La publicación rectangular recupera
\begin{equation}
 R_{\rm el}=\sqrt{b_S/a_S}=e^{-\eta/2},\qquad
 R_{\rm el}^4=\frac{1-\xi}{1+\xi},\qquad
 \xi=\frac{1-R_{\rm el}^4}{1+R_{\rm el}^4}.
 \label{ii:eq:ii-elipse-reciproca}
\end{equation}
La prueba consiste en sustituir
\(2\eta=\log((1+\xi)/(1-\xi))\). El intercambio de ejes
envía \(\eta\) a \(-\eta\), \(R_{\rm el}\) a
\(R_{\rm el}^{-1}\) y \(\tau=iR_{\rm el}^2\) a \(-1/\tau\).
Se conserva \(a_Sb_S\). El radio \(R_{\rm el}\) es adimensional;
el radio polar \(\rho_{\partial}\) tiene unidad de raíz de acción.

El interior admite además la distancia proyectiva de Hilbert. Si
\(z_-,x,y,z_+\) están ordenados en una cuerda de la elipse, se define
\[
 d_H(x,y)=\frac12\log
 \frac{|y-z_-|\,|x-z_+|}{|x-z_-|\,|y-z_+|}.
\]
Es una distancia adimensional: las unidades cancelan en la razón doble.
La normalización \((Q,P)\mapsto(Q/a_S,P/b_S)\) lleva la elipse
al disco unidad y preserva esa razón, pues las longitudes sobre una
misma recta se multiplican por un factor común. Transporta así la métrica
al modelo de Beltrami--Klein. Ni esta distancia interior ni la distancia
entre focos se identifica con una longitud atómica.

\subsection{Radio de Bohr: consecuencia de la escala de acción}
\label{ii:sec:ii-elipse-bohr}

En la realización electrodinámica, la escala de Bohr se expresa como
\(a_{\rm B}=\hbar/(m_ec\alpha)\) \cite{ii:codata2022}. En la cadena
del presente artículo, esa lectura sucede a las construcciones de
acción, vacío, estructura fina y masa electrónica. Se reciben aquí
\(\hbar,m_e,c,\alpha>0\) en una misma realización dimensional.
\begin{proposition}[Expresión areal de la escala de Bohr]
\label{ii:prop:ii-elipse-bohr}
En esa realización,
\begin{equation}
 a_{\rm B}=\frac{a_Sb_S}{2m_ec\alpha}
 =\frac{\mathcal A(2\pi)}{2\pi m_ec\alpha}
 =\frac{\bar\lambda_C(m_e)}{\alpha},\qquad
 \bar\lambda_C(m_e)=\frac{\hbar}{m_ec}.
 \label{ii:eq:ii-elipse-bohr}
\end{equation}
\end{proposition}
\begin{proof}
Las identidades \(a_Sb_S=2\hbar\) y
\(\mathcal A(2\pi)=2\pi\hbar\), sustituidas en la expresión
de Bohr, dan los dos primeros miembros. La definición de
\(\bar\lambda_C(m_e)\) da el tercero.
\end{proof}

El numerador tiene dimensión de acción y \(m_ec\) de momento; el
cociente es una longitud. Esta composición utiliza el producto de
semiejes y las otras tres magnitudes: no afirma que Bohr sea un semieje,
un foco, una distancia de Hilbert o el radio modular. Tampoco equipara
la masa electrónica con el cuanto circular de orden 108. La notación
\(a_{\rm B}\) evita confundir esta longitud con la normalización
\(a_0\) del operador de área.

\subsection{Realización inductiva--capacitiva y continuidad constitutiva}
\label{ii:sec:ii-elipse-lc}

El puente electrodinámico de la elipse conserva una carta adicional:
un reloj \(\omega>0\) y una impedancia de referencia \(Z_*>0\).
La construcción recuperada del corpus define
\begin{equation}
 L_\eta=\frac{Z_*}{\omega}e^{-\eta},\qquad
 C_\eta=\frac{1}{Z_*\omega}e^\eta,
 \qquad Z_\eta=\sqrt{L_\eta/C_\eta}.
 \label{ii:eq:ii-elipse-lc}
\end{equation}
Aquí \(C_\eta\) es capacitancia, distinta del contraángulo \(C^*\).
\begin{proposition}[Producto dinámico y cociente constitutivo]
\label{ii:prop:ii-elipse-lc}
En la carta anterior,
\begin{equation}
 L_\eta C_\eta=\omega^{-2},\qquad
 \frac{Z_\eta}{Z_*}=e^{-\eta}
 =\frac{b_S}{a_S}=R_{\rm el}^{2}.
 \label{ii:eq:ii-elipse-impedancia}
\end{equation}
Con \(Q=\sqrt{Z_*}\,q\), \(P=\Phi/\sqrt{Z_*}\), la energía
\(H_{LC}=q^2/(2C_\eta)+\Phi^2/(2L_\eta)\) toma la forma
\[
 H_{LC}=\frac\omega2(e^{-\eta}Q^2+e^\eta P^2).
\]
Sobre \(\gamma(\theta)\), sus componentes eléctrica y magnética son
\(U_E=\hbar\omega\cos^2\theta\) y
\(U_B=\hbar\omega\sin^2\theta\); su suma es \(\hbar\omega\).
\end{proposition}
\begin{proof}
El producto en \eqref{ii:eq:ii-elipse-lc} cancela los exponentes y
\(Z_*\); su cociente es \(Z_*^2e^{-2\eta}\). Tomar la raíz
positiva y usar \eqref{ii:eq:ii-elipse-reciproca} da
\eqref{ii:eq:ii-elipse-impedancia}. Sustituir \(q=Q/\sqrt{Z_*}\)
y \(\Phi=\sqrt{Z_*}P\) en \(H_{LC}\) da su forma equilibrada.
Finalmente, los semiejes de \(\gamma\) cancelan los factores
\(e^{\mp\eta}\), y \(\cos^2\theta+\sin^2\theta=1\).
\end{proof}

Este puente distingue la forma, el reloj y el transductor. La elección
del sentido dinámico debe concordar con la convención simpléctica: para
las ecuaciones \(\dot Q=\partial_PH\),
\(\dot P=-\partial_QH\), la parametrización anterior satisface
\(\dot\theta=-\omega\). Las identidades energéticas no dependen
de invertir ese recorrido.

El lector resolvente del vacío actúa sobre los canales ya producidos
\(q_\pm=e^{-(A_{\rm rad}\pm C^*_{\rm rad})}\) y la incidencia
\(90/120\). Sus respuestas
\(r_\pm=(1-q_\pm^{90})/(1-q_\pm^{120})\) publican
\(\widehat Z=r_-/r_+\) y \(\widehat c=(r_+r_-)^{-1}\).
Son respuestas espectrales, no radios geométricos. La identidad
\eqref{ii:eq:ii-elipse-impedancia} fija aquí la realización LC;
identificarla con la impedancia constitutiva del vacío conserva el mapa
entre ambos lectores y sus cartas. No se igualan sus valores por compartir
el antecedente angular.
